\begin{textsample}{POS Dim 1 – 2006 – Score 21.00 – 2006/t000160.txt}  \label{ex:f1_pos_012}
\textbf{The} future that we will create can be a future that we ' ll be proud of.
I think about this every day ; it ' s quite literally my job.
I ' m co-founder and senior columnist at Worldchanging. com.
Alex Steffen and I founded Worldchanging in late 2003, and since then we and our growing global team of contributors have documented \textbf{the} ever-expanding variety of solutions that are out there, right now and on \textbf{the} near horizon.
In a little over two years, we ' ve written up about 4, 000 items — replicable models, technological tools, emerging ideas — all providing a path to a future that ' s more sustainable, more equitable and more desirable.
Our emphasis on solutions is quite intentional.
There are tons of places to go, online and off, if what you want to find is \textbf{the} latest bit of news about just how quickly our hell-bound handbasket is moving.
We want to offer people an idea of what they can do about it.
We focus primarily on \textbf{the} \textbf{planet} ' s environment, but we also address issues of global development, international conflict, responsible use of emerging technologies, even \textbf{the} rise of \textbf{the} so-called Second Superpower and much, much more. \textbf{The} scope of solutions that we discuss is actually pretty broad, but that reflects both \textbf{the} range of challenges that need to be met and \textbf{the} kinds of innovations that will allow us to do so.
A quick sampling really can barely scratch \textbf{the} \textbf{surface}, but to give you a sense of what we cover : tools for rapid disaster relief, such as this inflatable concrete shelter ; innovative uses of bioscience, such as a flower that changes color in \textbf{the} presence of landmines ; ultra high-efficiency designs for homes and offices ; distributed power generation using solar power, wind power, ocean power, other clean energy sources ; ultra, ultra high-efficiency vehicles of \textbf{the} future ; ultra high- efficiency vehicles you can get right now ; and better urban design, so you don ' t need to drive as much in \textbf{the} first place ; bio-mimetic approaches to design that take advantage of \textbf{the} efficiencies of natural models in both vehicles and buildings ; distributed computing projects that will help us model \textbf{the} future of \textbf{the} climate.
Also, a number of \textbf{the} topics that we ' ve been talking about this week at TED are things that we ' ve addressed in \textbf{the} past on Worldchanging : cradle-to-cradle design, MIT ' s Fab Labs, \textbf{the} consequences of extreme longevity, \textbf{the} One \textbf{Laptop} per Child project, even Gapminder.
As a born-in-the-mid-1960s Gen X-er, hurtling all too quickly to my fortieth birthday, I ' m naturally inclined to pessimism.
But working at Worldchanging has convinced me, much to my own surprise, that successful responses to \textbf{the} world ' s problems are nonetheless possible.
Moreover, I ' ve come to realize that focusing only on negative outcomes can really blind you to \textbf{the} very possibility of success.
As Norwegian social scientist Evelin Lindner has observed,"Pessimism is a luxury of good times...
In difficult times, pessimism is a self- fulfilling, self-inflicted death sentence."\textbf{The} truth is, we can build a better world, and we can do so right now.
We have \textbf{the} tools : we saw a hint of that a moment ago, and we ' re coming up with new ones all \textbf{the} time.
We have \textbf{the} knowledge, and our understanding of \textbf{the} \textbf{planet} improves every day.
Most importantly, we have \textbf{the} motive : we have a world that needs fixing, and nobody ' s going to do it for us.
Many of \textbf{the} solutions that I and my colleagues seek out and write up every day have some important aspects in common : transparency, collaboration, a willingness to experiment, and an appreciation of science — or, more appropriately, science! \textbf{The} majority of models, tools and ideas on Worldchanging encompass combinations of these characteristics, so I want to give you a few concrete examples of how these principles combine in world- changing ways.
We can see world-changing values in \textbf{the} emergence of tools to make \textbf{the} invisible visible — that is, to make apparent \textbf{the} conditions of \textbf{the} world around us that would otherwise be largely imperceptible.
We know that people often change their behavior when they can see and understand \textbf{the} impact of their actions.
As a small example, many of us have experienced \textbf{the} change in driving behavior that comes from having a real time display of mileage showing precisely how one ' s driving habits affect \textbf{the} vehicle ' s efficiency. \textbf{The} last few years have all seen \textbf{the} rise of innovations in how we measure and display aspects of \textbf{the} world that can be too big, or too intangible, or too slippery to grasp easily.
Simple technologies, like wall-mounted \textbf{devices} that display how much power your household is using, and what kind of results you ' ll get if you turn off a few lights — these can actually have a direct positive impact on your energy footprint.
Community tools, like text messaging, that can tell you when pollen counts are up or smog levels are rising or a natural disaster is unfolding, can give you \textbf{the} information you need to act in a timely fashion.
Data-rich displays like maps of campaign contributions, or maps of \textbf{the} \textbf{disappearing} polar ice caps, allow us to better understand \textbf{the} context and \textbf{the} flow of processes that affect us all.
We can see world-changing values in research projects that seek to meet \textbf{the} world ' s medical needs through open access to data and collaborative action.
Now, some people emphasize \textbf{the} risks of knowledge-enabled dangers, but I ' m convinced that \textbf{the} benefits of knowledge-enabled solutions are far more important.
For example, open-access journals, like \textbf{the} Public Library of Science, make cutting-edge scientific research free to all — everyone in \textbf{the} world.
And actually, a growing number of science publishers are adopting this model.
Last year, hundreds of volunteer \textbf{biology} and \textbf{chemistry} researchers around \textbf{the} world worked together to sequence \textbf{the} \textbf{genome} of \textbf{the} parasite responsible for some of \textbf{the} developing world ' s worst diseases : African sleeping sickness, leishmaniasis and Chagas disease.
That \textbf{genome} data can now be found on open-access genetic data banks around \textbf{the} world, and it ' s an enormous boon to researchers trying to come up with treatments.
But my favorite example has to be \textbf{the} global response to \textbf{the} SARS epidemic in 2003, 2004, which relied on worldwide access to \textbf{the} full gene sequence of \textbf{the} SARS virus. \textbf{The} U.
S.
National Research Council in its follow-up report on \textbf{the} outbreak specifically cited this open availability of \textbf{the} sequence as a key reason why \textbf{the} treatment for SARS could be developed so quickly.
And we can see world-changing values in something as humble as a cell phone.
I can probably count on my fingers \textbf{the} number of people in this room who do not use a mobile phone — and where is Aubrey, because I know he doesn ' t?
For many of us, cell phones have really become almost an extension of ourselves, and we ' re really now beginning to see \textbf{the} social changes that mobile phones can bring about.
You may already know some of \textbf{the} big-picture aspects : globally, more camera phones were sold last year than any other kind of camera, and a growing number of people live lives mediated through \textbf{the} lens, and over \textbf{the} network — and sometimes enter history books.
In \textbf{the} developing world, mobile phones have become economic drivers.
A study last year showed a direct correlation between \textbf{the} growth of mobile phone use and subsequent GDP increases across Africa.
In Kenya, mobile phone minutes have actually become an alternative currency. \textbf{The} political aspects of mobile phones can ' t be ignored either, from text message swarms in Korea helping to bring down a government, to \textbf{the} Blairwatch Project in \textbf{the} UK, keeping tabs on politicians who try to avoid \textbf{the} press.
And it ' s just going to get more wild.
Pervasive, always-on networks, high quality sound and video, even \textbf{devices} made to be worn instead of carried in \textbf{the} \textbf{pocket}, will transform how we live on a scale that few really appreciate.
It ' s no exaggeration to say that \textbf{the} mobile phone may be among \textbf{the} world ' s most important technologies.
And in this rapidly \textbf{evolving} context, it ' s possible to imagine a world in which \textbf{the} mobile phone becomes something far more than a medium for social interaction.
I ' ve long admired \textbf{the} Witness project, and Peter Gabriel told us more details about it on Wednesday, in his profoundly moving presentation.
And I ' m just incredibly happy to see \textbf{the} news that Witness is going to be opening up a Web portal to enable users of digital cameras and camera phones to send in their recordings over \textbf{the} Internet, rather than just hand-carrying \textbf{the} videotape.
Not only does this add a new and potentially safer avenue for documenting abuses, it opens up \textbf{the} program to \textbf{the} growing global digital generation.
Now, imagine a similar model for networking environmentalists.
Imagine a Web portal collecting recordings and evidence of what ' s happening to \textbf{the} \textbf{planet} : putting news and data at \textbf{the} fingertips of people of all kinds, from activists and researchers to businesspeople and political figures.
It would highlight \textbf{the} changes that are underway, but would more importantly give voice to \textbf{the} people who are willing to work to see a new world, a better world, come about.
It would give everyday citizens a chance to play a role in \textbf{the} protection of \textbf{the} \textbf{planet}.
It would be, in essence, an"Earth Witness"project.
Now, just to be clear, in this talk I ' m using \textbf{the} name"Earth Witness"as part of \textbf{the} scenario, simply as a shorthand, for what this imaginary project could aspire to, not to piggyback on \textbf{the} wonderful work of \textbf{the} Witness organization.
It could just as easily be called,"Environmental Transparency Project,""Smart Mobs for Natural Security"— but Earth Witness is a lot easier to say.
Now, many of \textbf{the} people who participate in Earth Witness would focus on ecological problems, human-caused or otherwise, especially environmental crimes and significant sources of \textbf{greenhouse} gases and \textbf{emissions}.
That ' s understandable and important.
We need better documentation of what ' s happening to \textbf{the} \textbf{planet} if we ' re ever going to have a chance of repairing \textbf{the} damage.
But \textbf{the} Earth Witness project wouldn ' t need to be limited to problems.
In \textbf{the} best Worldchanging tradition, it might also serve as a showcase for good ideas, successful projects and efforts to make a difference that deserve much more visibility.
Earth Witness would show us two worlds : \textbf{the} world we ' re leaving behind, and \textbf{the} world we ' re building for generations to come.
And what makes this scenario particularly appealing to me is we could do it today. \textbf{The} key components are already widely available.
Camera phones, of course, would be fundamental to \textbf{the} project.
And for a lot of us, they ' re as close as we have yet to always-on, widely available information tools.
We may not remember to bring our digital cameras with us wherever we go, but very few of us forget our phones.
You could even imagine a version of this scenario in which people actually build their own phones.
Over \textbf{the} course of last year, \textbf{open-source} hardware hackers have come up with multiple models for usable, Linux-based mobile phones, and \textbf{the} Earth Phone could spin off from this kind of project.
At \textbf{the} other end of \textbf{the} network, there ' d be a server for people to send photos and messages to, accessible over \textbf{the} Web, combining a photo-sharing service, social networking platforms and a collaborative filtering system.
Now, you Web 2. 0 folks in \textbf{the} audience know what I ' m talking about, but for those of you for whom that last sentence was in a crazy moon language, I mean simply this : \textbf{the} online part of \textbf{the} Earth Witness project would be created by \textbf{the} users, working together and working openly.
That ' s enough right there to start to build a compelling chronicle of what ' s now happening to our \textbf{planet}, but we could do more.
An Earth Witness site could also serve as a \textbf{collection} spot for all sorts of data about conditions around \textbf{the} \textbf{planet} picked up by environmental sensors that attach to your cell phone.
Now, you don ' t see these \textbf{devices} as add-ons for phones yet, but students and engineers around \textbf{the} world have attached atmospheric sensors to bicycles and handheld \textbf{computers} and cheap robots and \textbf{the} backs of pigeons — that being a project that ' s actually underway right now at U.
C.
Irvine, using bird-mounted sensors as a way of measuring smog-forming pollution.
It ' s hardly a stretch to imagine putting \textbf{the} same thing on a phone carried by a person.
Now, \textbf{the} idea of connecting a sensor to your phone is not new : phone-makers around \textbf{the} world offer phones that sniff for bad breath, or tell you to worry about too much sun exposure.
Swedish firm Uppsala Biomedical, more seriously, makes a mobile phone add-on that can process \textbf{blood} tests in \textbf{the} field, uploading \textbf{the} data, displaying \textbf{the} results.
Even \textbf{the} Lawrence Livermore National Labs have gotten into \textbf{the} act, designing a prototype phone that has radiation sensors to find dirty bombs.
Now, there ' s an enormous variety of tiny, inexpensive sensors on \textbf{the} market, and you can easily imagine someone putting together a phone that could measure temperature, CO2 or methane levels, \textbf{the} presence of some biotoxins — potentially, in a few years, maybe even H5N1 avian flu virus.
You could see that some kind of system like this would actually be a really good fit with Larry Brilliant ' s InSTEDD project.
Now, all of this data could be \textbf{tagged} with geographic information and mashed up with online maps for easy viewing and analysis.
And that ' s worth noting in particular. \textbf{The} impact of open-access online maps over \textbf{the} last year or two has been simply phenomenal.
Developers around \textbf{the} world have come up with an amazing variety of ways to layer useful data on top of \textbf{the} maps, from bus routes and crime statistics to \textbf{the} global progress of avian flu.
Earth Witness would take this further, linking what you see with what thousands or millions of other people see around \textbf{the} world.
It ' s kind of exciting to think about what might be accomplished if something like this ever existed.
We ' d have a far better — far better knowledge of what ' s happening on our \textbf{planet} environmentally than could be gathered with satellites and a handful of government sensor nets alone.
It would be a collaborative, bottom-up approach to environmental awareness and protection, able to respond to emerging concerns in a smart mobs kind of way — and if you need greater sensor density, just have more people show up.
And most important, you can ' t ignore how important mobile phones are to global youth.
This is a system that could put \textbf{the} next generation at \textbf{the} front lines of gathering environmental data.
And as we work to figure out ways to mitigate \textbf{the} worst effects of climate disruption, every little bit of information matters.
A system like Earth Witness would be a tool for all of us to participate in \textbf{the} improvement of our knowledge and, ultimately, \textbf{the} improvement of \textbf{the} \textbf{planet} itself.
Now, as I suggested at \textbf{the} outset, there are thousands upon thousands of good ideas out there, so why have I spent \textbf{the} bulk of my time telling you about something that doesn ' t exist?
Because this is what tomorrow could look like : bottom-up, technology-enabled global collaboration to handle \textbf{the} biggest crisis our civilization has ever faced.
We can save \textbf{the} \textbf{planet}, but we can ' t do it alone — we need each other.
Nobody ' s going to fix \textbf{the} world for us, but working together, making use of technological innovations and human communities alike, we might just be able to fix it ourselves.
We have at our fingertips a cornucopia of compelling models, powerful tools, and innovative ideas that can make a meaningful difference in our \textbf{planet} ' s future.
We don ' t need to wait for a magic \textbf{bullet} to save us all ; we already have an arsenal of solutions just waiting to be used.
There ' s a staggering \textbf{array} of wonders out there, across diverse disciplines, all telling us \textbf{the} same thing : success can be ours if we ' re willing to try.
And as we say at Worldchanging, another world isn ' t just possible ; another world is here.
We just need to open our eyes.
Thank you very much.

% matched lemmas: array, biology, blood, bullet, chemistry, collection, computer, device, disappear, emission, evolve, genome, greenhouse, laptop, open-source, planet, pocket, surface, tag, the
\end{textsample}
